\chapter{The complete runnable source}
\label{app:source}
These listings use byte-for-byte copies of the runnable files in the supplied project. The preparation script refreshes them automatically before compilation. They are the same code explained in Chapters~5 and~10. The line numbers are for reading; they are not part of the programs. The numerical results and figures elsewhere in the guide were generated by running these files.

\section{The model: src/MoranModel.jl}
\lstinputlisting[firstnumber=1]{code/MoranModel.jl}

\section{The demonstration: scripts/demo.jl}
\lstinputlisting[firstnumber=1]{code/demo.jl}

\section{The full study script}
\lstinputlisting[firstnumber=1]{code/run_study.jl}

\section{The automated checks: test/runtests.jl}
The Test library marks expressions that must be true. \code{@test} checks an ordinary assertion. \code{@test_throws} checks that an invalid request raises the expected error rather than returning a misleading number. \code{@testset} groups related checks. The function \code{isapprox} tests numerical agreement within tolerance; exact equality would be inappropriate for independently rounded floating-point calculations. The Monte Carlo check has a deliberately generous tolerance measured in standard errors; it does not demand that all nominal 95\% confidence intervals cover.
\lstinputlisting[firstnumber=1]{code/runtests.jl}

\section{Preparing the guide's plot inputs}
This helper does not add new biological rules or new stochastic experiments. It reads the completed full study, separates rows by population size, calculates error-bar lengths, and writes small data files that LaTeX can plot. It also prepares the printed tables and summary counts. \code{readlines} reads a file as text lines; \code{split} separates a CSV row at commas; \code{parse} converts a text field to a number. The dot in \code{String.(...)} applies conversion to each field. \code{push!} appends a row to an array.

The first preparation block also copies the runnable source and trajectory data into the guide folder. \code{cp} copies a file, and \code{force=true} allows replacing an older generated copy. \code{basename} keeps just the filename. \code{readdir} lists the data directory; \code{startswith} and \code{endswith} identify the trajectory CSV files. These copies keep the LaTeX project self-contained without asking its compiler to read outside its folder. Edit the original runnable code, then run this helper again; do not edit the generated copies separately.

\code{mkpath} creates the output directory if needed. \code{open(..., "w") do file} writes through the file handle and closes it at the end. The helper checks for 28 conditions and 20,000 repetitions per condition so that a quick run is not accidentally used as the full-study report. An inequality added to \code{covered} contributes one for true and zero for false. The standardised discrepancy divides the estimate's error by the exact binomial sampling standard error. This differs from the plug-in standard error stored in the study CSV, which uses the observed proportion.

The first loop writes summary numbers. The next writes the selected \(N=100\) table. For each population size, later loops prepare estimated points, a smooth exact curve, and the approximation-error data. The final loop prepares the initial-frequency comparison. The \code{@printf} instructions format LaTeX table rows; doubled backslashes in the string represent literal backslashes required by LaTeX. They do not alter a scientific calculation.

\lstinputlisting[firstnumber=1,basicstyle=\ttfamily\fontsize{8.5}{10.2}\selectfont]{code/prepare_guide.jl}
